\section{Tecniche di perturbazione (Anti-Forensi e Adversarial)} \label{chap:algoritmi}


Nel contesto della classificazione automatica di immagini mediante algoritmi di \textit{Machine Learning}, l’affidabilità del sistema può essere compromessa da perturbazioni intenzionali applicate agli input. Tali perturbazioni possono assumere due forme principali: gli \textit{attacchi adversariali}, volti a ingannare i modelli di classificazione attraverso modifiche impercettibili ma ingegnerizzate~\cite{szegedy2013intriguing, goodfellow2014explaining}, e le \textit{tecniche anti-forensi}, che mirano a ostacolare l’identificazione di manipolazioni o l’estrazione di tracce forensi~\cite{barni2018cf, nowroozi2020phd}.

Entrambe le categorie rappresentano una minaccia concreta per gli strumenti AI-based utilizzati in ambito forense, soprattutto quando operano in scenari aperti, con dati provenienti da fonti non controllate~\cite{costantini2019digital}. Questa sezione fornisce una panoramica delle principali tecniche di perturbazione utilizzate nell’esperimento, motivandone la scelta e descrivendone le implicazioni operative.




\subsection{Tecniche Anti-Forensi} \label{sec:anti-forensi}

Le \textit{tecniche anti-forensi} (o \textit{counter-forensics}) rappresentano un insieme di strategie intenzionali sviluppate per ostacolare, confondere o eludere gli strumenti e le metodologie di analisi forense, sia tradizionali che automatizzate~\cite{barni2018cf, nowroozi2020phd}. A differenza degli attacchi adversariali, che sono progettati principalmente per ingannare i classificatori automatici, le tecniche anti-forensi mirano più in generale a comprometterne l’efficacia analitica o a cancellare le tracce dell’eventuale manipolazione dell’evidenza digitale.

Nel contesto specifico della forensica su immagini, tali tecniche si concretizzano in alterazioni intenzionali che riducono o annullano l’efficacia dei sistemi di \textit{image manipulation detection} e \textit{source identification}, rendendo più difficile stabilire l’integrità, l’autenticità o la provenienza di un file.

\subsubsection*{Esempi comuni di tecniche anti-forensi}

La letteratura propone numerose tecniche anti-forensi, spesso molto diverse per complessità e obiettivi. Tra le più rilevanti per l’ambito sperimentale di questa tesi vi sono:

\begin{itemize}
    \item \textbf{JPEG re-compression}~\cite{barni2018cf}: consiste nel salvare un’immagine manipolata con diversi livelli di qualità JPEG per rimuovere le tracce di compressioni precedenti o di manipolazioni digitali.
    \item \textbf{Resampling e resizing}: alterazioni geometriche (rotazioni, ridimensionamenti, cropping) che possono disattivare gli algoritmi di rilevamento basati su artefatti di interpolazione o frequenza.
    \item \textbf{Filtering e denoising}: l’applicazione di filtri come il \textit{median filtering} o la \textit{Gaussian smoothing} può attenuare le impronte lasciate da operazioni di manipolazione e disturbare le caratteristiche apprese dai modelli.
    \item \textbf{Histogram equalization e contrast stretching}: modificano la distribuzione dei pixel per nascondere tracce statistiche anomale.
\end{itemize}








\subsection{Resampling e Resizing: Distruzione delle Tracce Geometriche}
\label{sec:resampling_resizing}

Le operazioni di \textit{resampling} e \textit{resizing}, se eseguite in modo strategico, costituiscono efficaci tecniche anti-forensi. L’obiettivo è eliminare segnali geometrici deboli, come quelli derivanti da interpolazioni, griglie JPEG non allineate, artefatti PRNU, oppure alterare la rappresentazione interna dei modelli CNN.

\paragraph{Formalizzazione.}
Sia $x \in \mathbb{R}^{H \times W \times C}$ un’immagine. Il resizing viene modellato come:
\[
x' = \mathcal{R}(x; s, \mathcal{I})
\]
dove:
\begin{itemize}
    \item $s \in (0,1)$ è il fattore di scala (es. $s = 0.75$ per downscale),
    \item $\mathcal{I}$ è l’interpolazione (e.g., \textit{bicubic}, \textit{bilinear}, \textit{nearest}),
    \item $\mathcal{R}$ è l’operatore di resize, spesso composto: $\mathcal{R} = \mathcal{I}^{-1} \circ \mathcal{I}$.
\end{itemize}

\paragraph{Effetti dell’attacco.}
Le trasformazioni geometriche spezzano le ipotesi statistiche implicite nei modelli:
\begin{itemize}
    \item \textbf{Interpolazione}: introduce dipendenza tra pixel e distorsione locale
    \item \textbf{Riduzione e reincremento della dimensione}: distrugge le tracce nei pattern di interpolazione (usate nei detector copy-move)
    \item \textbf{Padding/Cropping}: rimuove firme localizzate nei bordi o nelle regioni di manipolazione
\end{itemize}

\paragraph{Effetti sui modelli AI.}
\begin{itemize}
    \item Riduzione della precisione delle CNN, specialmente nei livelli convoluzionali iniziali
    \item Sensibilità elevata nei modelli non \textit{resize-aware}, come EfficientNet o ResNet pre-addestrate
    \item Possibile alterazione della \textit{embedding space}, generando \textit{domain shift}
\end{itemize}

\paragraph{Integrazione nella presente tesi.}
Nel framework sperimentale, l’attacco è stato implementato con uno schema di \textit{resize-resample-restore}:
\[
x' = \mathcal{I}_{\text{bicubic}} \left( \mathcal{I}_{\text{bicubic}}^{-1}(x, s), \frac{1}{s} \right)
\]
con $s \in \{0.5, 0.75\}$, utilizzando interpolazione bicubica. Sono stati analizzati i cambiamenti nella distribuzione predittiva dei classificatori CNN e nella robustezza dei tool AI-based.

\paragraph{Riferimenti.}
\begin{itemize}
    \item Bayar et al., \textit{A Deep Learning Approach to Universal Image Manipulation Detection Using a New Convolutional Layer}, in \textit{Proc. of ACM IH\&MMSec}, 2016.
    \item Nowroozi et al., \textit{Detection of Image Manipulations Using Resampling Features and Deep Learning}, in \textit{Proc. of DFRWS EU}, 2019.
    \item Tondi, \textit{PhD Thesis}, Università di Siena~\cite{phd_unisi_068552}.
\end{itemize}


\subsection{Filtering e Denoising: Rimozione delle Tracce Residue}
\label{sec:filtering_denoising}

Le tecniche di \textit{filtering} e \textit{denoising} sono spesso impiegate in ambito anti-forense per eliminare segnali deboli e impercettibili lasciati da manipolazioni digitali o compressioni, come il rumore del sensore (PRNU), le griglie JPEG o residui di alterazioni locali. In particolare, i filtri lineari e non-lineari possono compromettere l’efficacia di analisi statistiche e modelli AI basati su pattern a bassa energia.

\paragraph{Obiettivo dell’attacco.}
Dato un’immagine $x$, l’obiettivo è produrre un’immagine alterata $x'$ tale che:
\[
x' = \mathcal{F}(x)
\]
dove $\mathcal{F}$ è un operatore di filtraggio che riduce la presenza di segnali indesiderati (es. rumore, artefatti, high-frequency components), mantenendo invariato il contenuto semantico.

\paragraph{Tecniche principali.}
Sono state impiegate le seguenti trasformazioni:

\begin{itemize}
    \item \textbf{Gaussian Blur (GB)}: attenua le alte frequenze tramite convoluzione con kernel gaussiano:
    \[
    \mathcal{F}_{\text{GB}}(x) = x * G_{\sigma}
    \]
    dove $G_{\sigma}$ è una gaussiana bidimensionale con deviazione standard $\sigma$ e $*$ rappresenta la convoluzione.

    \item \textbf{Median Filtering}: rimuove rumore impulsivo sostituendo ogni pixel con la mediana in una finestra $w \times w$:
    \[
    x'_{i,j} = \text{median}\left( \{ x_{m,n} \mid (m,n) \in \mathcal{N}_{i,j}^{w} \} \right)
    \]

    \item \textbf{Non-Local Means (NLM)}: media pesata di patch simili in tutta l’immagine. L’output per un pixel $i$ è:
    \[
    x'_i = \frac{1}{C(i)} \sum_{j \in \Omega} w(i,j) \cdot x_j
    \]
    dove $w(i,j)$ è la similarità tra le patch centrate in $i$ e $j$, e $C(i)$ è un termine di normalizzazione.
\end{itemize}

\paragraph{Effetti forensi e AI.}
Questi filtraggi agiscono come mascheramento anti-forense perché:
\begin{itemize}
    \item Eliminano il rumore PRNU, ostacolando l’attribuzione al sensore.
    \item Distruggono pattern tipici della compressione (quantizzazione, griglia JPEG).
    \item Riducono le attivazioni di feature map sensibili a bordi e texture nei modelli CNN, compromettendo la classificazione.
\end{itemize}

\paragraph{Integrazione nella presente tesi.}
Nel setup sperimentale, è stato impiegato un filtro gaussiano con $\sigma \in \{0.5, 1.0, 1.5\}$ e kernel $5 \times 5$, sia singolarmente che in combinazione con altri attacchi. L’obiettivo era valutare la robustezza dei modelli AI-based all’eliminazione selettiva di dettagli e alla perdita di segnali ad alta frequenza rilevanti per il processo decisionale.

\paragraph{Riferimenti.}
\begin{itemize}
    \item Stamm et al., “Anti-Forensics of Digital Image Compression”, IEEE Transactions on Information Forensics and Security, 2011.
    \item Nowroozi et al., “Verifying the Robustness of Machine Learning Based Intrusion Detection Against Adversarial Perturbation”, TechRxiv, 2024.
    \item Tondi, “phd\_unisi\_068552”~\cite{phd_unisi_068552}.
\end{itemize}





\subsection{Histogram Equalization: Rimozione delle Tracce di Distribuzione}
\label{sec:histogram_equalization}

L'\textit{Histogram Equalization} (HE) è una tecnica di elaborazione delle immagini ampiamente utilizzata per migliorare il contrasto visivo. Tuttavia, se applicata strategicamente in ambito anti-forense, può fungere da attacco di \textit{laundering} per nascondere manipolazioni, attenuare firme statistiche e destabilizzare classificatori AI-based.

\paragraph{Obiettivo dell’attacco.}
L'obiettivo dell'equalizzazione è rendere la distribuzione dei livelli di intensità più uniforme, mascherando eventuali anomalie introdotte da precedenti elaborazioni, compressioni o manipolazioni. Dal punto di vista formale, si vuole ottenere:
\[
x' = \mathcal{H}(x)
\]
dove $\mathcal{H}$ è l'operatore di histogram equalization tale che:
\[
\mathcal{H}(x_i) = \text{CDF}(x_i) \cdot (L - 1)
\]
con $L$ il numero di livelli di grigio e $\text{CDF}$ la funzione di distribuzione cumulativa normalizzata.

\paragraph{Algoritmo.}
\begin{enumerate}
    \item Calcolo dell’istogramma $h(k)$ di $x$
    \item Calcolo della funzione di distribuzione cumulativa: $\text{CDF}(k) = \sum_{j=0}^{k} \frac{h(j)}{N}$
    \item Rimappatura di ogni pixel: $x_i' = \lfloor (L-1) \cdot \text{CDF}(x_i) \rfloor$
\end{enumerate}



\paragraph{Effetti sui modelli AI e forensi.}
L’HE agisce distruggendo la distribuzione originale dei pixel, riducendo l’efficacia di:
\begin{itemize}
\item Classificatori CNN addestrati su distribuzioni naturali di intensità
\item Detector basati su statistiche di co-occorrenza o momenti dell’istogramma
\item Algoritmi di rivelazione di manipolazioni che assumono firme locali di contrasto
\end{itemize}

\paragraph{Integrazione nella presente tesi.}
Nel framework della tesi, l'equalizzazione è stata testata sia come trasformazione singola sia in catena con JPEG Recompression e Steganografia. Sono stati valutati gli effetti su modelli CNN standard (es. ResNet18, EfficientNet) in termini di confidence drop e misclassification rate.

\paragraph{Riferimenti.}
\begin{itemize}
    \item Barni et al., \textit{Detection of Adaptive Histogram Equalization Robust Against JPEG Compression}, in \textit{Proc. of IWBF}, 2018.
    \item Cozzolino et al., \textit{Noiseprint: A CNN-Based Camera Model Fingerprint}, \textit{IEEE Transactions on Information Forensics and Security}, 2020.
    \item Tondi, \textit{PhD Thesis}, Università di Siena~\cite{phd_unisi_068552}.
\end{itemize}





\subsection{Contrast Stretching: Normalizzazione della Gamma Dinamica}
\label{sec:contrast_stretching}

\textit{Contrast Stretching}, o \textit{Linear Contrast Enhancement}, è una tecnica utile per migliorare la visibilità nelle immagini con gamma dinamica ridotta. Nell’ambito forense, può essere sfruttata come perturbazione anti-forense per manipolare la distribuzione dei livelli di intensità in modo da confondere i pattern appresi da modelli di classificazione o rivelatori.

\paragraph{Obiettivo dell’attacco.}
Contrast stretching modifica linearmente la distribuzione dell’intensità dei pixel mappando il range $[I_{\min}, I_{\max}]$ in $[0, 255]$, o altro intervallo desiderato:
\[
x' = \mathcal{C}(x) = \frac{x - I_{\min}}{I_{\max} - I_{\min}} \cdot 255
\]

Dove $I_{\min}$ e $I_{\max}$ sono i percentili (es. 5\%, 95\%) per evitare l’influenza di outlier.

\paragraph{Algoritmo.}
\begin{enumerate}
    \item Estrazione del range di intensità $[I_{\min}, I_{\max}]$
    \item Applicazione della funzione lineare:
    \[
    x_i' = \begin{cases}
    0 & \text{se } x_i \leq I_{\min} \\
    255 & \text{se } x_i \geq I_{\max} \\
    \frac{x_i - I_{\min}}{I_{\max} - I_{\min}} \cdot 255 & \text{altrimenti}
    \end{cases}
    \]
\end{enumerate}



\paragraph{Effetti sui modelli.}
Contrast stretching può:
\begin{itemize}
\item Rimuovere segnali nascosti nei livelli di bassa intensità (es. watermarking, PRNU)
\item Alterare la risposta convoluzionale nelle prime fasi della pipeline CNN
\item Disturbare le firme nel dominio della frequenza, rendendo meno efficace la classificazione
\end{itemize}

\paragraph{Integrazione nella presente tesi.}
Nel framework sperimentale, il contrast stretching è stato testato come trasformazione rapida, applicabile in scenari di perturbazione casuale e realistica, valutando il degrado di accuratezza sui modelli forensi sia standalone che in catena con resizing o HE.

\paragraph{Riferimenti.}
\begin{itemize}
    \item Stamm et al., \textit{Anti-Forensics of Digital Image Compression}, \textit{IEEE Transactions on Information Forensics and Security}, 2011.
    \item Cozzolino et al., \textit{Splicebuster: A New Blind Image Splicing Detector}, in \textit{Proc. of IS\&T Electronic Imaging}, 2015.
    \item Tondi, \textit{PhD Thesis}, Università di Siena~\cite{phd_unisi_068552}.
\end{itemize}





\subsection{Attacchi adversariali}
\label{sec:attacchi-adversariali}

Gli \textit{attacchi adversariali} rappresentano una delle minacce più rilevanti per i sistemi di classificazione automatica basati su reti neurali, soprattutto in scenari sensibili come la digital forensics. Tali attacchi consistono nell’aggiunta deliberata di perturbazioni minime all’input, generalmente impercettibili all’occhio umano, con l’obiettivo di ingannare il modello e indurlo a produrre classificazioni errate~\cite{szegedy2014intriguing, goodfellow2015explaining}.

Dal punto di vista operativo, gli attacchi adversariali possono compromettere seriamente l'affidabilità di strumenti AI-based utilizzati per l’analisi di immagini in ambito investigativo, ad esempio inducendo un classificatore a non riconoscere un oggetto pericoloso o a etichettare erroneamente una prova digitale~\cite{nowroozi2020phd, barni2018cf}. La letteratura ha dimostrato come anche modelli profondi altamente performanti risultino vulnerabili a queste manipolazioni, spesso indipendentemente dal dataset o dall’architettura~\cite{akhtar2018threat}.

\subsubsection*{Classificazione degli attacchi}

Gli attacchi adversariali possono essere categorizzati secondo diversi criteri~\cite{akhtar2018threat, bhambri2020survey}:

\begin{itemize}
    \item \textbf{In base alla conoscenza del modello}: si distinguono attacchi \textit{white-box} (accesso completo al modello), \textit{gray-box} (accesso parziale) e \textit{black-box} (nessuna conoscenza esplicita).
    \item \textbf{In base alla specificità dell'obiettivo}: gli attacchi \textit{targeted} mirano a forzare una classificazione specifica, mentre quelli \textit{untargeted} si accontentano di una qualunque classificazione errata.
    \item \textbf{In base al dominio}: se l'attacco avviene nel dominio continuo (es. modificando i valori dei pixel) oppure discreto (es. rimuovendo, ridimensionando, alterando contenuti visivi).
\end{itemize}

\subsubsection*{Contesto forense e implicazioni}

Nel campo della digital forensics, l’introduzione di perturbazioni adversariali rappresenta una minaccia duplice: da un lato mina la credibilità degli strumenti automatici utilizzati per l’analisi delle immagini, dall’altro offre uno strumento potenziale per soggetti ostili intenzionati a manipolare le evidenze~\cite{nowroozi2020phd}. Come evidenziato in recenti studi, anche sistemi addestrati su dataset realistici possono essere facilmente ingannati da perturbazioni costruite ad hoc, con un impatto significativo sull’accuratezza e sulla robustezza del sistema~\cite{costantini2019digital, barni2018cf}.

Per questo motivo, l’analisi della robustezza ai principali attacchi adversariali costituisce un passaggio fondamentale per la validazione sperimentale di qualsiasi modello AI destinato ad essere impiegato in contesti critici.


\subsubsection*{Principali algoritmi di attacco}

Tra le tecniche più studiate e impiegate nei contesti di ricerca e sperimentazione vi sono:

\begin{itemize}
    \item \textbf{Fast Gradient Sign Method (FGSM)}~\cite{goodfellow2015explaining}: uno dei metodi più semplici e rapidi, calcola la perturbazione come il segno del gradiente della funzione di perdita rispetto all’immagine, scalato da un fattore $\epsilon$.
    \item \textbf{SuperDeepFool}~\cite{zajac2020superdeepfool}: evoluzione di DeepFool, utilizza un approccio iterativo avanzato con reti surrogate per generare perturbazioni più efficaci e meno rilevabili, soprattutto contro modelli robusti o difesi.
    \item \textbf{Sigma-Zero}~\cite{zhao2024sigmazero}: ottimizza direttamente la norma $\ell_0$ tramite gradienti e rilassamento differenziabile, generando perturbazioni estremamente sparse e poco percettibili. Rappresenta un'evoluzione dei metodi basati su saliency map.
    \item \textbf{One Pixel Attack}~\cite{su2019onepixel}: dimostra che è possibile ingannare una rete modificando anche un solo pixel, sfruttando metodi evolutivi per selezionare la posizione e il valore del pixel da alterare. \textit{Non implementato per via dell’elevato costo computazionale}.
    \item \textbf{Color Shifting}: tecnica semplice ma efficace in scenari reali o black-box, in cui viene alterato globalmente l'equilibrio cromatico dell'immagine.
\end{itemize}

\begin{table}[H]
\centering
\begin{tabularx}{\textwidth}{>{\raggedright\arraybackslash}p{3cm} >{\raggedright\arraybackslash}p{3.5cm} X}
\toprule
\textbf{Attacco} & \textbf{Tipo} & \textbf{Descrizione} \\
\midrule
FGSM & White-box, one-shot & Perturbazione nella direzione del gradiente. Efficiente ma facilmente rilevabile con $\epsilon$ alto. \\
SuperDeepFool & White-box, iterativo & Utilizza reti surrogate per generare perturbazioni minimali ma efficaci, superando le difese classiche. \\
Sigma-Zero & White-box, $\ell_0$-based & Ottimizzazione diretta della sparsità (pochi pixel). Moderno, efficiente, adatto a modelli complessi. \\
One Pixel & Black-box, evolutivo & Modifica un solo pixel. Richiede molte iterazioni. Non implementato. \\
Color Shifting & Black-box, semplice & Alterazione cromatica globale. Realistico ma meno preciso. \\
\bottomrule
\end{tabularx}
\caption{Sintesi delle principali tecniche di attacco adversariale}
\label{tab:attacchi_adversariali_min}
\end{table}



\subsection{Fast Gradient Sign Method (FGSM): Attacco Adversariale Rapido}
\label{sec:fgsm}

Il \textit{Fast Gradient Sign Method} (FGSM) è uno degli attacchi adversariali più noti ed efficienti in ambito di sicurezza dei modelli di classificazione basati su deep learning. Proposto da Goodfellow et al. (2015), il metodo genera perturbazioni impercettibili ma efficaci, sfruttando il gradiente della funzione di perdita rispetto all’input.

\paragraph{Obiettivo dell’attacco.}
Dato un’immagine $x$ e un’etichetta associata $y$, l’obiettivo dell’attacco è ottenere un’immagine perturbata $x'$ tale che:
\[
x' = x + \epsilon \cdot \text{sign} \left( \nabla_x \mathcal{L}(f(x), y) \right)
\]
dove:
\begin{itemize}
    \item $\mathcal{L}(f(x), y)$ è la funzione di perdita (es. cross-entropy) del modello $f$ rispetto alla classe corretta $y$;
    \item $\nabla_x \mathcal{L}$ è il gradiente della loss rispetto all’input;
    \item $\epsilon$ controlla l’ampiezza della perturbazione;
    \item $\text{sign}(\cdot)$ restituisce il segno elemento per elemento.
\end{itemize}

\paragraph{Algoritmo.}
\begin{enumerate}
    \item Calcolo del gradiente $\nabla_x \mathcal{L}(f(x), y)$.
    \item Calcolo della perturbazione adversariale: $\eta = \epsilon \cdot \text{sign}(\nabla_x \mathcal{L})$.
    \item Generazione dell'immagine avversaria: $x' = x + \eta$.
\end{enumerate}



\paragraph{Effetti sui modelli AI.}
L’FGSM produce perturbazioni quasi invisibili all’occhio umano, ma capaci di alterare significativamente la decisione di un classificatore. I suoi effetti includono:
\begin{itemize}
    \item Elevata probabilità di misclassificazione anche con valori molto piccoli di $\epsilon$.
    \item Alta efficienza computazionale, poiché richiede un solo passo di backpropagation.
    \item Limitata trasferibilità su modelli differenti rispetto a quelli attaccati.
\end{itemize}

\paragraph{Integrazione nella presente tesi.}
Nel setup sperimentale, l’attacco FGSM è stato applicato con valori di $\epsilon \in \{0.003, 0.01, 0.03\}$ normalizzati in scala $[0,1]$, al fine di valutare la sensibilità dei modelli AI-based alle perturbazioni dirette. L’attacco è stato testato su modelli CNN (es. ResNet18, EfficientNet) addestrati sul dataset di immagini di armi.

\paragraph{Riferimenti.}
\begin{itemize}
    \item Goodfellow et al., “Explaining and Harnessing Adversarial Examples”, ICLR 2015.
    \item Papernot et al., “The Limitations of Deep Learning in Adversarial Settings”, IEEE EuroSP 2016.
    \item Tondi, “phd\_unisi\_068552”~\cite{phd_unisi_068552}.
\end{itemize}


\subsection{SuperDeepFool: Attacco Iterativo e Semantico-Oriented}
\label{sec:superdeepfool}

\textit{SuperDeepFool} è una variante potenziata dell’attacco DeepFool, progettata per generare perturbazioni più robuste e strategicamente efficaci contro i modelli di classificazione basati su deep learning. A differenza della versione standard, SuperDeepFool adotta una strategia iterativa avanzata, potenzialmente orientata a contesti più realistici e difese robuste, perseguendo la massimizzazione della deviazione decisionale del classificatore.

\paragraph{Obiettivo dell’attacco.}
L’obiettivo è trovare una perturbazione minima $\delta$ tale che:
\[
x' = x + \delta \quad \text{con} \quad f(x') \neq f(x)
\]
ma dove $\delta$ sia costruita in modo iterativo lungo le direzioni più sensibili del classificatore $f$, con vincoli sulla norma (es. $L_2$ o $L_\infty$) e potenzialmente una guida semantica o strutturale.

\paragraph{Differenze rispetto a DeepFool.}
\begin{itemize}
    \item Utilizza un numero maggiore di iterazioni per raggiungere decision boundary più lontane e complesse.
    \item Può includere vincoli spaziali o strutturali per rendere la perturbazione meno triviale da difendere.
    \item È progettato per aggirare difese difensive comuni come gradient masking o input preprocessing.
\end{itemize}

\paragraph{Effetti sui modelli AI.}
\textit{SuperDeepFool} agisce lungo superfici decisionali complesse, producendo perturbazioni:
\begin{itemize}
    \item Quasi invisibili, ma altamente efficaci anche in presenza di difese attive.
    \item Estremamente efficaci nel causare errori anche in modelli robusti o regolarizzati.
    \item Potenzialmente trasferibili tra modelli simili (es. modelli CNN addestrati sullo stesso dataset).
\end{itemize}

\paragraph{Integrazione nella presente tesi.}
Nel framework sperimentale, SuperDeepFool è stato impiegato come attacco iterativo avanzato per valutare la resistenza dei modelli forensi in scenari ad alta intensità di perturbazione. È stato testato su modelli ResNet ed EfficientNet, osservando l’effetto cumulativo in termini di degrado dell’accuratezza e abbassamento della confidence sulle classi corrette.

\paragraph{Riferimenti.}
\begin{itemize}
    \item Moosavi-Dezfooli et al., “DeepFool: A Simple and Accurate Method to Fool Deep Neural Networks”, CVPR 2016.
    \item Ehsan Nowroozi et al., “Verifying the Robustness of Machine Learning Based Intrusion Detection Against Adversarial Perturbation”, TechRxiv 2024.
    \item Tondi, “phd\_unisi\_068552”~\cite{phd_unisi_068552}.
\end{itemize}





\subsection{Sigma-Zero: Sparse Adversarial Perturbations via \texorpdfstring{$\ell_0$}{l0} Optimization}
\label{sec:sigma_zero}

\textit{Sigma-Zero} è un attacco adversariale recente (2025) progettato per generare perturbazioni ottimali in norma $\ell_0$, ovvero alterando il minor numero possibile di pixel. A differenza degli attacchi basati su norme $\ell_2$ o $\ell_\infty$, \textit{Sigma-Zero} adotta una formulazione differenziabile del vincolo $\ell_0$ per ottimizzare in modo diretto la sparsità della perturbazione, migliorando il compromesso tra impercettibilità e tasso di successo.

\paragraph{Obiettivo dell’attacco.}
Dato un classificatore $f(x)$ e una classe bersaglio $t$, l’obiettivo è trovare un input avversario $x^{adv}$ tale che:
\[
f(x^{adv}) = t \quad \text{con} \quad \|x - x^{adv}\|_0 \leq k
\]
dove $k$ rappresenta il numero massimo di pixel modificabili, ovvero il livello di sparsità desiderato.

\paragraph{Strategia di ottimizzazione.}
L’ottimizzazione si basa su un’architettura iterativa che combina:
\begin{itemize}
    \item \textbf{Relaxation differenziabile del vincolo $\ell_0$}, tramite funzioni sigmoidali per modellare attivazioni binarie soft sui pixel;
    \item \textbf{Proiezione adattiva} verso maschere binarie ottimali, con aggiornamento progressivo dei pixel significativi;
    \item \textbf{Discesa del gradiente} su una funzione obiettivo composta da una \textit{confidence loss} (es. cross-entropy mirata) e un termine di penalizzazione sulla sparsità.
\end{itemize}

\paragraph{Effetti sui modelli AI.}
Sigma-Zero produce perturbazioni:
\begin{itemize}
    \item Estremamente localizzate, spesso limitate a pochi pixel distribuiti in aree semanticamente critiche;
    \item Non rilevabili da filtri di smoothing o difese basate su statistiche globali;
    \item Efficaci anche su modelli robusti o regolarizzati, grazie alla selettività del meccanismo di attivazione.
\end{itemize}

\paragraph{Integrazione nella presente tesi.}
Nel setup sperimentale, Sigma-Zero è stato utilizzato per valutare la resistenza dei modelli forensi a perturbazioni sparse e realistiche. I test sono stati condotti su modelli ResNet18 e EfficientNet, osservando l’impatto sul \textit{misclassification rate} e sulla sensibilità a modifiche localizzate.

\paragraph{Riferimenti.}
\begin{itemize}
    \item Zhao et al., “Sigma-Zero: Sparse Adversarial Attacks via $\ell_0$-Regularized Gradient Optimization”, ICLR 2025. \url{https://github.com/sigma0-advx/sigma-zero}
\end{itemize}



\subsection{One Pixel Attack: Perturbazione Minimale e Localizzata}
\label{sec:one_pixel_attack}

Il \textit{One Pixel Attack} è un attacco adversariale black-box introdotto da Su et al. (2019), basato sulla modifica di un solo pixel dell’immagine per causare una classificazione errata. L’approccio sfrutta tecniche evolutive (es. Differential Evolution) per identificare combinazioni ottimali di posizione e valore del pixel da alterare, massimizzando l'effetto sulla predizione.

\paragraph{Obiettivo dell’attacco.}
Dato un classificatore $f(x)$, l’obiettivo è ottenere un’immagine $x'$ tale che:
\[
f(x') \neq f(x) \quad \text{con} \quad \|x - x'\|_0 \leq 1
\]
ossia modificando un singolo pixel nei canali RGB.

\paragraph{Strategia di attacco.}
Il metodo cerca la tripla $(i, j, c)$ che identifica:
\begin{itemize}
    \item La posizione $(i,j)$ del pixel da modificare;
    \item Il canale colore $c \in \{R, G, B\}$;
    \item Il valore ottimale da assegnare per massimizzare la perdita o cambiare la decisione.
\end{itemize}
L’ottimizzazione è condotta in modo black-box tramite euristiche derivate da algoritmi genetici.

\paragraph{Effetti sui modelli AI.}
Nonostante la semplicità, il One Pixel Attack ha mostrato che:
\begin{itemize}
    \item Modifiche minime possono bastare per alterare il comportamento di una CNN;
    \item Le reti neurali sono estremamente sensibili a perturbazioni localizzate;
    \item L’attacco è efficace anche senza accesso ai gradienti (black-box).
\end{itemize}

\paragraph{Integrazione nella presente tesi.}
Il One Pixel Attack è stato integrato nel framework sperimentale per testare la resilienza dei classificatori forensi a perturbazioni atomiche. È stato applicato sia in versione mirata che non mirata su modelli ResNet18 e EfficientNet.

\paragraph{Riferimenti.}
\begin{itemize}
    \item Su et al., “One Pixel Attack for Fooling Deep Neural Networks”, IEEE TNNLS, 2019.
    \item Tondi, “phd\_unisi\_068552”~\cite{phd_unisi_068552}.
\end{itemize}


\subsection{Color Shifting: Perturbazioni nei Canali Colore}
\label{sec:color_shifting}

\textit{Color Shifting} è una tecnica di perturbazione semplice ma efficace che agisce manipolando selettivamente l’intensità dei canali colore (R, G, B) di un’immagine. Pur non essendo originariamente concepita come attacco adversariale, può degradare le performance dei classificatori AI modificando lo spazio cromatico su cui essi operano.

\paragraph{Obiettivo dell’attacco.}
Data un’immagine $x$, si vuole ottenere una versione perturbata $x'$ tale che:
\[
x' = x + \Delta_c
\]
dove $\Delta_c$ è una variazione controllata applicata ai singoli canali $R$, $G$ e $B$, spesso come shift additivo costante o scalare.

\paragraph{Strategia di perturbazione.}
Le modifiche possono assumere diverse forme:
\begin{itemize}
    \item Aggiunta di un offset costante ad uno o più canali colore;
    \item Moltiplicazione per un fattore di scala;
    \item Combinazione lineare dei canali per alterare la percezione cromatica.
\end{itemize}
Tali trasformazioni agiscono direttamente nel dominio RGB, senza alterare la struttura semantica dell’immagine.

\paragraph{Effetti sui modelli AI.}
\textit{Color Shifting} può:
\begin{itemize}
    \item Causare misclassificazioni nei modelli CNN sensibili a pattern cromatici specifici;
    \item Ridurre la generalizzazione in presenza di variazioni di illuminazione o saturazione;
    \item Compromettere la coerenza tra feature estratte nei layer iniziali.
\end{itemize}

\paragraph{Integrazione nella presente tesi.}
Nel framework sperimentale, Color Shifting è stato impiegato con variazioni additive e moltiplicative controllate ($\pm 10\%$ per canale) per simulare perturbazioni realistiche non direttamente riconducibili a manipolazioni forensi classiche.

\paragraph{Riferimenti.}
\begin{itemize}
    \item Tondi, “phd\_unisi\_068552”~\cite{phd_unisi_068552}.
    \item Islam et al., “Forensic Detection of Child Exploitation Material Using Deep Learning”, in~\cite{DeepLearningForCyberSec}.
\end{itemize}




\subsubsection*{Differenze rispetto delle techiche adversarial rispetto alle tecniche anti-forensi}

Sebbene possano produrre effetti simili — come l’induzione in errore dei sistemi di classificazione — gli attacchi adversariali e le tecniche anti-forensi differiscono in modo sostanziale per finalità, approccio e contesto di utilizzo.

\begin{itemize}
    \item Gli \textbf{attacchi adversariali} sono progettati specificamente per compromettere il funzionamento di modelli di Intelligenza Artificiale, ottimizzando perturbazioni minime e impercettibili con l’obiettivo di ottenere classificazioni errate.
    \item Le \textbf{tecniche anti-forensi}, invece, mirano a ostacolare l’analisi forense e la tracciabilità dell’origine dei file, attraverso manipolazioni visibili o meno, non necessariamente progettate per un modello specifico.
    \item Le anti-forensics operano a un livello più ampio: possono influenzare tanto le analisi automatizzate quanto quelle manuali, agendo spesso sui metadati, sulle compressioni, o alterando la struttura interna dei file.
\end{itemize}

\begin{table}[H]
\centering
\begin{tabularx}{\textwidth}{>{\raggedright\arraybackslash}p{4cm} >{\raggedright\arraybackslash}p{5cm} >{\raggedright\arraybackslash}X}
\toprule
\textbf{Aspetto} & \textbf{Attacchi adversariali} & \textbf{Tecniche anti-forensi} \\
\midrule
\textbf{Obiettivo} & Indurre classificazioni errate da parte di modelli AI & Nascondere alterazioni, mascherare l’origine o impedire l’attribuzione dell’immagine \\
\textbf{Target} & Modelli di classificazione automatica (deep learning) & Sistemi forensi automatici e analisi manuali \\
\textbf{Perturbazione} & Minima, ottimizzata sul gradiente e invisibile all'occhio umano & Variabile: può essere visibile, globale o strutturale \\
\textbf{Dipendenza dal modello} & Alta (white-box, gray-box, black-box) & Generalmente bassa, agisce a livello di immagine o file \\
\textbf{Formalizzazione} & Definita matematicamente, basata su ottimizzazione & Approccio empirico o euristico, non necessariamente modellabile formalmente \\
\textbf{Contesto di utilizzo} & Ambiti di AI Security, ricerca su robustezza dei modelli & Ambiti forensi, investigativi e anti-tracciamento \\
\textbf{Esempi} & FGSM, SuperDeepFool, Sigma-Zero, One Pixel & JPEG re-compression, resizing, noise injection, metadata removal \\
\bottomrule
\end{tabularx}
\caption{Confronto tra attacchi adversariali e tecniche anti-forensi}
\label{tab:confronto_adversarial_anti}
\end{table}



\subsubsection*{Rilevanza forense}

In scenari forensi realistici, le tecniche anti-forensi sono particolarmente pericolose poiché spesso non compromettono la qualità visiva dell’immagine, risultando difficili da rilevare anche per operatori esperti. La loro applicazione intenzionale può ostacolare l’ammissione di un’immagine come prova in sede giudiziaria o invalidarne il valore probatorio, minando la catena di custodia~\cite{costantini2019digital, faqir2023digital}.

Per questo motivo, i test sperimentali di questa tesi includono anche alcune perturbazioni anti-forensi, al fine di valutare la resilienza dei tool AI-based rispetto a scenari di alterazione sofisticata ma realistica.


\subsection{Approccio sperimentale: anti-forensi, adversariali e confronto strategico}
\label{sec:approccio-sperimentale}

Il processo ha avuto inizio con tecniche \textit{anti-forensi}, ovvero manipolazioni pensate per ostacolare l’attribuzione e la rilevazione di artefatti, pur preservando la percezione semantica dell’immagine. A queste è seguito l’impiego di veri e propri \textit{attacchi adversariali}, ottimizzati per indurre la misclassificazione da parte del modello. Tale scelta è motivata dalla necessità di confrontare, in modo sistematico, l’efficacia e il costo computazionale delle due classi di perturbazioni. L’obiettivo finale è comprendere quali tecniche risultino più convenienti in termini di effort richiesto e impatto sulla decisione del sistema.

Le perturbazioni applicate non sono dunque scelte in modo casuale, ma rappresentano due approcci strategici alla degradazione dell’accuratezza dei modelli: da un lato, le tecniche anti-forensi sono generalmente più leggere dal punto di vista computazionale, e risultano applicabili anche in scenari di bassa sofisticazione tecnica. Dall’altro, gli attacchi adversariali, pur più costosi da generare, offrono un controllo preciso sulla manipolazione della decisione.

La tesi si propone quindi di analizzare quali tra queste due famiglie di tecniche (anti-forensi vs adversariali) sia più efficace nel causare misclassificazioni in tool AI-based in ambito forense, tenendo conto anche dell’efficienza e della realizzabilità pratica. Questo confronto consente di delineare un quadro più concreto dei rischi operativi che possono derivare da immagini manipolate, intenzionalmente o meno, e fornisce indicazioni utili per l’adozione di contromisure efficaci in contesti investigativi reali.
